\chapter{Filing an Issue}\label{cha:issues}

This machine has one author, and an issue you file will, in all
likelihood, be handed straight to an AI coding session to work on, as
written. That is worth knowing, because it changes what a good report
is: not a description of the fault, but the machine's own account of it.
The machine can write most of that account itself, and the whole of this
appendix is how to let it.

\section{Before you try again}
None of these is politeness. Each one removes a day.

\textbf{Turn \texttt{DUMP ON} before you reproduce it.} Then every
fifteen seconds, and again the instant you type \verb|DUMP| once it has
gone wrong, the machine writes its whole state --- registers, the screen
as you see it, the last 4096 instructions, your keystrokes, and the log
of every command you have typed --- to \pth{dumps/dump-NNN.txt} on the
host. A dump taken while it is wrong is worth more than any description
either of us could write, and \verb|DUMP ON| means you cannot miss the
moment.

\textbf{Try it once with \texttt{k4510 {-}{-}no-startup.bat}.} A
\pth{/STARTUP.BAT} that has gone wrong looks exactly like a machine that
has gone wrong, and telling them apart has cost this project a day
before. Ten seconds of yours settles it.

\textbf{Check this book.} If the machine disagrees with a page here, that
is still a fault worth filing --- it is this book's, and it is fixed the
same way. Say which page.

\section{BUG writes the report}
Then let the machine do the writing:

\begin{type}
/] BUG
\end{type}

\noindent \verb|BUG| interviews you --- seven questions, one at a time,
each answered in a line or two --- and writes the finished report to

\begin{center}
\texttt{/SYSTEM/LOG/BUGREPORTS/BUG-YYYYMMDD-HHMMSS.TXT}
\end{center}

\noindent making the directory if it is not there. Five lines it fills
in itself, and they are the five most often got wrong by hand: which
machine (the K4510 on its own Linux, or in a window on somebody's
desktop), the exact build
(\verb|INFO -v| says the same --- a trailing \verb|+| means the emulator
was built from an edited tree, which is exactly the report that will not
reproduce), the screen you are in, the last dump, and when.

The questions are the ones a coding session would otherwise have to ask
you: where you were (the shell, a BASIC, CP/M\ldots), what you typed
\emph{exactly}, what happened, what you expected instead, and \emph{why}
you expected it --- the line most reports leave out. The K4510 is a
fantasy machine (page~\pageref{cha:disclaimer}): there is no silicon it
can be measured against, so ``wrong'' only means something next to what
you were expecting and where that came from. It is also what separates a
bug from a design decision from an error in this book, which are three
different repairs.

\verb|BUG| is a program rather than a command, so it reaches you
everywhere the \texttt{*} escape does: \verb|*BUG| from either BASIC,
from the monitor, from CP/M.

\begin{aside}
\textbf{\texttt{DUMP} first, then \texttt{BUG}.} If the report's
\texttt{Dump} line says there is none, the more valuable half is missing.
Type \verb|DUMP| while it is still wrong and run \verb|BUG| again ---
the two travel together, and this is the one thing that is easy to do in
the wrong order.
\end{aside}

\section{Where it goes}
\begin{center}
\url{https://github.com/mlongval/k4510/issues}
\end{center}

\noindent From the host, the report is in \pth{fs/SYSTEM/LOG/BUGREPORTS/} and
the dump in \pth{dumps/}: open an issue, attach both, and you are done ---
the report is the whole text of the issue. The form there asks
\verb|BUG|'s questions again for anyone who has no report file: a fault
in this book, a machine that will not boot, a telephone. If you have no
account there, the same two files in an email are worth exactly as much.

Feature requests and disagreements are welcome in the same shape ---
``what I expected'' does the work whether or not anything is broken.

\begin{aside}
\textbf{One thing to check before you attach a dump.} It carries the
shell log, which is every command line you typed in that session, and
file names from your own disk along with them. It is a plain text file:
open it and have a look before it goes anywhere public.
\end{aside}
